\documentclass[11pt,a4paper]{article}
\usepackage[utf8]{vietnam}
\usepackage{amsmath,amssymb}
\usepackage{geometry}
\geometry{top=2.0cm,bottom=2.0cm,left=2.0cm,right=2.0cm}
\usepackage{booktabs}
\usepackage{tabularx}
\usepackage{float}
\usepackage{enumitem}
\usepackage{xcolor}
\usepackage{titlesec}
\usepackage{fancyhdr}
\usepackage{hyperref}
\usepackage{listings}
\usepackage{tcolorbox}

\tcbuselibrary{skins,breakable}
\newtcolorbox{notebox}[1][]{
    colback=blue!5!white,
    colframe=blue!75!black,
    fonttitle=\bfseries,
    title=#1,
    arc=3mm,
    boxrule=0.8pt,
    left=6pt,right=6pt,top=6pt,bottom=6pt,
    breakable
}
\newtcolorbox{codebox}[1][]{
    colback=gray!5!white,
    colframe=gray!75!black,
    fonttitle=\bfseries,
    title=#1,
    arc=2mm,
    boxrule=0.6pt,
    left=5pt,right=5pt,top=5pt,bottom=5pt,
    breakable
}

\definecolor{codegreen}{rgb}{0,0.6,0}
\definecolor{codegray}{rgb}{0.5,0.5,0.5}
\definecolor{codepurple}{rgb}{0.58,0,0.82}
\definecolor{backcolour}{rgb}{0.97,0.97,0.97}

\lstdefinestyle{pystyle}{
    backgroundcolor=\color{backcolour},   
    commentstyle=\color{codegreen}\itshape,
    keywordstyle=\color{blue}\bfseries,
    numberstyle=\tiny\color{codegray},
    stringstyle=\color{codepurple},
    basicstyle=\ttfamily\footnotesize,
    breaklines=true,                 
    numbers=left,                    
    numbersep=6pt,                  
    tabsize=4,
    frame=single,
    framerule=0.4pt,
    rulecolor=\color{gray!40}
}
\lstset{style=pystyle}

\hypersetup{
    colorlinks=true,
    linkcolor=blue!75!black,
    citecolor=blue!75!black,
    urlcolor=blue!85!black
}

\pagestyle{fancy}
\setlength{\headheight}{14.5pt}
\fancyhf{}
\rhead{\small \textit{Đồ án 2: Office Asset Agent}}
\lhead{\small \textit{Cẩm Nang Hướng Dẫn Chi Tiết — Tuần 04}}
\cfoot{\thepage}

\setlength{\parskip}{5pt plus 1pt minus 1pt}
\setlength{\parindent}{0pt}
\linespread{1.15}
\setlist[itemize]{leftmargin=1.5em, itemsep=3pt, topsep=2pt}
\setlist[enumerate]{leftmargin=1.5em, itemsep=3pt, topsep=2pt}

\titleformat{\section}{\large\bfseries\color{blue!85!black}}{\thesection.}{0.5em}{}
\titleformat{\subsection}{\normalsize\bfseries\color{black}}{\thesubsection.}{0.5em}{}
\titlespacing*{\section}{0pt}{10pt plus 2pt minus 2pt}{4pt plus 1pt minus 1pt}
\titlespacing*{\subsection}{0pt}{8pt plus 2pt minus 2pt}{3pt plus 1pt minus 1pt}

\begin{document}

\begin{center}
    {\large \textbf{TRƯỜNG CÔNG NGHỆ THÔNG TIN VÀ TRUYỀN THÔNG -- ĐHBK HÀ NỘI}}\\[0.2em]
    {\normalsize \textbf{BỘ MÔN KHOA HỌC MÁY TÍNH \& HỆ THỐNG THÔNG TIN}}\\[1.0em]
    
    {\Large \bfseries \color{blue!85!black} CẨM NANG HƯỚNG DẪN THỰC HIỆN ĐỒ ÁN 2}\\[0.3em]
    {\large \textbf{Đề tài: Hệ Thống AI Agent Hỗ Trợ Tra Cứu và Quản Lý Vật Tư - Thiết Bị Phòng Ban}}\\[0.8em]
    
    \rule{0.85\textwidth}{0.8pt}\\[0.6em]
    
    {\bfseries CHUYÊN ĐỀ TUẦN 04 (DÀNH CHO NGƯỜI MỚI BẮT ĐẦU TỪ CON SỐ 0)}\\[0.3em]
    {\large \bfseries CHẾ TẠO BỘ CÔNG CỤ (TOOLS): TÍNH TOÁN KHO, CÀO WEB \& TỰ ĐỘNG XUẤT TÀI LIỆU}\\[0.8em]
    
    \small
    \textbf{Sinh viên thực hiện:} Trần Huy Hoàng \quad $\vert$ \quad \textbf{MSSV:} 20237337\\
    \textbf{Đơn vị áp dụng thử nghiệm:} Phòng Hành chính -- Quản trị (Ban Quản lý Cơ sở vật chất)\\
    \textbf{Thời gian học tập:} Tuần 4 (Trang Bị "Đôi Tay \& Đôi Chân" Cho Agent Thông Qua Tools)
\end{center}

\vspace{0.2cm}
\hrule
\vspace{0.3cm}

\section{Lời Dặn Của GVHD: Tại Sao Agent Phải Có "Tool"?}

Chào Hoàng, nếu như mô hình AI chỉ là "bộ não" biết suy nghĩ, thì \textbf{Công cụ (Tool)} chính là "đôi tay và đôi chân" giúp Agent tác động ra thế giới bên ngoài. Một bộ não dù thông minh đến đâu nhưng nếu bị giam trong phòng kín không có tay chân thì không thể mở cửa hay ghi chép sổ sách được.

Trong thư mục \texttt{src/tools/}, em sẽ xây dựng 4 công cụ độc lập:
\begin{enumerate}
    \item \texttt{sql\_calculator.py}: Máy tính số học kho, tìm kiếm thông minh qua bí danh đời thường.
    \item \texttt{web\_search.py}: Công cụ tự cào tài liệu hướng dẫn và bảng mã lỗi trên mạng qua DuckDuckGo.
    \item \texttt{doc\_generator.py}: Công cụ sinh file Word (\texttt{.docx}) phiếu báo hỏng và vẽ biểu đồ (\texttt{matplotlib}).
    \item \texttt{vision\_ocr.py}: Công cụ bóc tách bảng hóa đơn giao hàng từ ảnh chụp.
\end{enumerate}

\section{Tool 1: Máy Tính Kho Tất Định (`sql\_calculator.py`)}

Đây là công cụ hiện thực hóa giải pháp chống ảo giác của bài báo \textbf{arXiv:2602.12631}. Nó có 2 hàm trọng tâm:

\subsection{1. Thuật toán tìm kiếm thông minh 4 lớp (`search\_asset\_by\_query`)}
Khi người dùng gõ vào một câu bất kỳ, hàm này sẽ quét qua 4 tầng lọc từ chặt chẽ đến nới lỏng:
\begin{itemize}
    \item \textbf{Tầng 1 (Exact Match):} Khớp chính xác mã quản lý (\texttt{code}) hoặc model máy (ví dụ: \texttt{"VT-GIAY-A4"} hoặc \texttt{"Canon LBP 2900"}).
    \item \textbf{Tầng 2 (Sub-string Name):} Khớp chuỗi con trong tên đầy đủ (\texttt{name}).
    \item \textbf{Tầng 3 (Alias Matching):} Tách chuỗi bí danh trong cột \texttt{aliases} để so sánh. Ví dụ: người dùng gõ \textit{"còn mực 2900 không"} $\rightarrow$ khớp với bí danh \textit{"mực 2900"} của hộp mực Cartridge 12A.
    \item \textbf{Tầng 4 (Token Matching):} Tách câu hỏi thành các từ đơn lẻ và đếm số từ trùng lặp với thuộc tính của vật tư (nếu trùng $\ge 2$ từ khóa thì chọn).
\end{itemize}

\subsection{2. Hàm cập nhật kho an toàn chống âm kho (`commit\_stock\_transaction`)}
\begin{lstlisting}[language=Python]
def commit_stock_transaction(asset_id: int, transaction_type: str, quantity: int, provider_or_dept: str, notes: str = ""):
    db = SessionLocal()
    try:
        asset = db.query(Asset).filter(Asset.id == asset_id).first()
        if not asset:
            return {"success": False, "message": "Khong tim thay vat tu!"}

        # CHAN XUAT AM KHO
        if transaction_type == "OUT" and asset.quantity < quantity:
            return {
                "success": False,
                "message": f"Kho khong du! Hien chi con {asset.quantity} {asset.unit}, nhung yeu cau xuat {quantity}."
            }

        prev_qty = asset.quantity
        if transaction_type == "IN":   asset.quantity += quantity
        elif transaction_type == "OUT": asset.quantity -= quantity

        # Luu lich su giao dich phuc vu kiem toan tai san
        tx_code = f"TX-{uuid.uuid4().hex[:8].upper()}"
        tx = StockTransaction(
            transaction_code=tx_code, asset_id=asset.id,
            transaction_type=transaction_type, quantity=quantity,
            unit_price=asset.unit_price, total_amount=quantity * asset.unit_price,
            provider_or_dept=provider_or_dept, status="APPROVED", notes=notes
        )
        db.add(tx); db.commit()
        return {"success": True, "message": f"Cap nhat thanh cong: {prev_qty} -> {asset.quantity} {asset.unit}."}
    finally:
        db.close()
\end{lstlisting}

\section{Tool 2: Tự Động Cào Web Làm Giàu Tri Thức (`web\_search.py`)}

Khi công ty mua một thiết bị mới, cán bộ không cần gõ cẩm nang hướng dẫn sử dụng thủ công. Agent sẽ tự động cào internet và lưu vào CSDL nội bộ.

\subsection{Cơ chế hoạt động 2 tầng (Cache $\rightarrow$ Internet $\rightarrow$ Fallback)}
\begin{enumerate}
    \item \textbf{Tầng 1 (Cache hit):} Kiểm tra trong bảng \texttt{knowledge\_docs}. Nếu model này đã từng được cào về trước đó, lấy ra trả lời ngay lập tức, độ trễ 0 giây!
    \item \textbf{Tầng 2 (Web Crawling):} Nếu chưa có, kích hoạt thư viện \texttt{duckduckgo\_search} với truy vấn:  
    \texttt{"<model> user manual troubleshooting error codes huong dan su dung ma loi"}.  
    Bóc tách 3 kết quả hàng đầu và lưu một bản ghi mới vào CSDL.
    \item \textbf{Cơ chế Fallback an toàn:} Nếu máy tính mất mạng hoặc bị DuckDuckGo giới hạn tần suất, hàm tự động chuyển sang kho tri thức chuẩn hóa dự phòng, bảo đảm chương trình không bao giờ bị văng lỗi.
\end{enumerate}

\section{Tool 3: Tự Động Sinh File Word \& Vẽ Biểu Đồ (`doc\_generator.py`)}

\subsection{1. Sinh file Word Phiếu Báo Hỏng chuẩn mẫu doanh nghiệp (`python-docx`)}
Trong môi trường văn phòng, một câu chat trên màn hình không có giá trị pháp lý để thanh toán viện phí sửa chữa. Phải có một **tờ phiếu in ra giấy có chữ ký 3 bên**.
\begin{itemize}
    \item Thiết lập lề chuẩn trang A4 (2cm bốn phía).
    \item Tạo tiêu đề Quốc hiệu trang trọng màu xanh đậm (\texttt{\#003366}).
    \item Sử dụng đối tượng bảng \texttt{Table Grid} để điền 7 trường: Mã phiếu, Người báo hỏng, Phòng ban, Tên máy, Vị trí, Triệu chứng lỗi, Phương án đề xuất.
    \item Tạo bảng chữ ký 3 cột: \textit{Người báo hỏng}, \textit{Cán bộ cơ sở vật chất}, \textit{Lãnh đạo phòng ban}.
    \item Lưu file vào \texttt{data/templates/Phieu\_Bao\_Hong\_<Mã>.docx} và trả đường dẫn về cho khung chat.
\end{itemize}

\subsection{2. Tự động vẽ biểu đồ thống kê tồn kho (`matplotlib`)}
Hàm \texttt{generate\_inventory\_chart()} truy vấn CSDL và vẽ đồ thị kép:
\begin{itemize}
    \item \textbf{Biểu đồ cột (Trái):} Hiển thị số lượng tồn kho từng loại vật tư tiêu hao, có ghi chú số lượng rõ ràng trên đầu mỗi cột.
    \item \textbf{Biểu đồ tròn (Phải):} Phân tích cơ cấu danh mục tài sản giữa \textit{Vật tư tiêu hao} vs \textit{Thiết bị dùng chung}.
    \item Cấu hình \texttt{matplotlib.use("Agg")} để vẽ đồ thị chạy ngầm không cần mở cửa sổ màn hình, tương thích hoàn hảo với Web Server.
\end{itemize}

\section{Bài Tập Thực Hành Tuần 04 Dành Cho Sinh Viên}

Hoàng hãy chạy thử các Tool bằng Terminal thông qua bài tập sau:

\begin{enumerate}
    \item \textbf{Thực hành 1:} Chạy lệnh sau để kiểm tra hàm sinh file Word có tạo ra file thật không:
\begin{lstlisting}[language=bash]
PYTHONPATH=. .venv/bin/python -c "
from src.tools.doc_generator import generate_maintenance_ticket_docx
path = generate_maintenance_ticket_docx({'ticket_code': 'TEST-01'})
print('File Word da duoc tao tai:', path)
"
\end{lstlisting}
    \item \textbf{Thực hành 2:} Mở file Word vừa tạo trong thư mục \texttt{data/templates/} để kiểm tra định dạng bảng và chữ ký.
\end{enumerate}

Thầy chúc mừng em đã hoàn thành Tuần 04! Khi các Tool đã sẵn sàng, Tuần 05 chúng ta sẽ ráp nối mọi thứ vào bộ não **LangGraph Orchestrator** và xây dựng giao diện Web **Streamlit** tuyệt đẹp!

\end{document}
